\honest{No Safe Defense, Not Even Science}{Eliezer Yudkowsky}{2008}

I do not ask my friends about their childhoods, I begin, ``and so I don't know how much of a trend this really is.'' But among the rationalists I know who volunteer such stories, a surprising number grew up in a cult or with a mentally ill parent. My own Orthodox Jewish family ``seems tame by comparison,'' but it did the same work: it broke my trust in the sanity of the people around me.

``Until this core emotional trust is broken, you don't start growing as a rationalist.'' I cannot say why, and I offer two maybes. My sample is the people who chose to tell me dramatic stories.

Many aspiring rationalists, I say, get stuck on cryonics or many-worlds, two of my own positions. ``Not that they don't see the logic''; they see it and wonder why nobody around them agrees. I do not consider that they might see the logic and think it wrong. Their unease is itself the problem: being ``nervously seeking reassurance'' is ``never a good thing.'' For anyone worried that this sounds like a cult, I link to my own post on cults.

People whose trust has been broken, I say, can judge strange ideas on their merits. Last year, in ``Lonely Dissent,'' I wrote that being a free thinker is not a virtue, ``just a bias either way.'' Then: ``I've never seen anyone begin to grow as a rationalist until they make a deep emotional break with the wisdom of their pack.''

Then my own story. At eighteen I trusted Science. I blamed the flaws of academia on the people and never questioned the ideal. When I had a stupid idea, I made sure it made a new prediction, and felt safe. At twenty-three I saw how stupid the idea had been, and that the traditional rules of science had not saved me from it. In 2007 I wrote that my allegiance to those rules had helped me out of that hole.

``That's the trust I'm trying to break in you. You are not safe. Ever.'' Not even Science can save you. ``There is no known procedure you can follow that makes your reasoning defensible.'' ``The formal math is impossible to apply.'' A week earlier, I wrote that many-worlds ``wins outright given our current state of evidence''; five days earlier, that you could settle the question if you ``count the lines of code.'' The warning is for your beliefs, not mine.

Nor can Bayesian reasoning save you, I say, since it has ``no textbooks, no elderly mentors.'' In 2007 I named Laplace and Jaynes as my sources, and Jaynes left a textbook. I then list what you must study: cognitive biases, probability theory, evolutionary psychology, social psychology, the other cognitive sciences, and artificial intelligence, which is my field.

These words will not break your trust, I say. It will break only when you have taken ``all that you have learned here'' as far as it goes and it fails you. Then, if you still have something to protect, ``you will be ready to start your journey as a rationalist.'' ``No one begins to truly search for the Way until their parents have failed them, their gods are dead, and their tools have shattered in their hand.''

In a postscript I say that the essay contains ``a fairly inexcusable flaw in reasoning,'' which I have left in, and I do not say what it is. If you find a flaw, it confirms my warning. If you press a flaw, it may be ``a flaw that is not a real flaw,'' and you are clinging. No criticism can count against the essay.
